\documentclass[10pt,a4paper]{amsart}
\usepackage{amsmath,amssymb,amsthm,graphicx}
\usepackage{verbatim}
\newenvironment{reference}{\comment}{\endcomment}
%\newenvironment{reference}{}{}
\newenvironment{slogan}{\comment}{\endcomment}
\newenvironment{history}{\comment}{\endcomment}

% For commutative diagrams we use Xy-pic
\usepackage[all]{xy}

% We use 2cell for 2-commutative diagrams.
\xyoption{2cell}
% Theorem environments.
%
\theoremstyle{plain}
\newtheorem{theorem}[subsection]{Theorem}
\newtheorem{proposition}[subsection]{Proposition}
\newtheorem{lemma}[subsection]{Lemma}

\theoremstyle{definition}
\newtheorem{definition}[subsection]{Definition}
\newtheorem{example}[subsection]{Example}
\newtheorem{exercise}[subsection]{Exercise}
\newtheorem{situation}[subsection]{Situation}

\theoremstyle{remark}
\newtheorem{remark}[subsection]{Remark}
\newtheorem{remarks}[subsection]{Remarks}

\numberwithin{equation}{subsection}

% Macros
%
\def\lim{\mathop{\mathrm{lim}}\nolimits}
\def\colim{\mathop{\mathrm{colim}}\nolimits}
\def\Spec{\mathop{\mathrm{Spec}}}
\def\Hom{\mathop{\mathrm{Hom}}\nolimits}
\def\Ext{\mathop{\mathrm{Ext}}\nolimits}
\def\SheafHom{\mathop{\mathcal{H}\!\mathit{om}}\nolimits}
\def\SheafExt{\mathop{\mathcal{E}\!\mathit{xt}}\nolimits}
\def\Sch{\mathit{Sch}}
\def\Mor{\mathop{\mathrm{Mor}}\nolimits}
\def\Ob{\mathop{\mathrm{Ob}}\nolimits}
\def\Sh{\mathop{\mathit{Sh}}\nolimits}
\def\NL{\mathop{N\!L}\nolimits}
\def\CH{\mathop{\mathrm{CH}}\nolimits}
\def\proetale{{pro\text{-}\acute{e}tale}}
\def\etale{{\acute{e}tale}}
\def\QCoh{\mathit{QCoh}}
\def\Ker{\mathop{\mathrm{Ker}}}
\def\Im{\mathop{\mathrm{Im}}}
\def\Coker{\mathop{\mathrm{Coker}}}
\def\Coim{\mathop{\mathrm{Coim}}}

% Boxtimes
%
\DeclareMathSymbol{\boxtimes}{\mathbin}{AMSa}{"02}

%
% Macros for moduli stacks/spaces
%
\def\QCohstack{\mathcal{QC}\!\mathit{oh}}
\def\Cohstack{\mathcal{C}\!\mathit{oh}}
\def\Spacesstack{\mathcal{S}\!\mathit{paces}}
\def\Quotfunctor{\mathrm{Quot}}
\def\Hilbfunctor{\mathrm{Hilb}}
\def\Curvesstack{\mathcal{C}\!\mathit{urves}}
\def\Polarizedstack{\mathcal{P}\!\mathit{olarized}}
\def\Complexesstack{\mathcal{C}\!\mathit{omplexes}}
% \Pic is the operator that assigns to X its picard group, usage \Pic(X)
% \Picardstack_{X/B} denotes the Picard stack of X over B
% \Picardfunctor_{X/B} denotes the Picard functor of X over B
\def\Pic{\mathop{\mathrm{Pic}}\nolimits}
\def\Picardstack{\mathcal{P}\!\mathit{ic}}
\def\Picardfunctor{\mathrm{Pic}}
\def\Deformationcategory{\mathcal{D}\!\mathit{ef}}
\usepackage{cleveref}
\crefname{originallemma}{original lemma}{original lemmas}
\Crefname{originallemma}{Original lemma}{Original lemmas}
\makeatletter
\newlabel{algebra-lemma-completion-generalities}{{0315}{source}{}{sourcepointer-1}{}}
\newlabel{algebra-lemma-completion-generalities@cref}{{[originallemma][0][]0315}{[1][0][]source}{}{}{}}
\makeatother
\crefname{originallemma}{original lemma}{original lemmas}
\Crefname{originallemma}{Original lemma}{Original lemmas}
\makeatletter
\newlabel{algebra-lemma-mod-injective}{{00ME}{source}{}{sourcepointer-2}{}}
\newlabel{algebra-lemma-mod-injective@cref}{{[originallemma][0][]00ME}{[1][0][]source}{}{}{}}
\makeatother
\crefname{originalreference}{original reference}{original references}
\Crefname{originalreference}{Original reference}{Original references}
\makeatletter
\newlabel{lemma-split-node}{{0CBW}{source}{}{sourcepointer-3}{}}
\newlabel{lemma-split-node@cref}{{[originalreference][0][]0CBW}{[1][0][]source}{}{}{}}
\makeatother
\crefname{originallemma}{original lemma}{original lemmas}
\Crefname{originallemma}{Original lemma}{Original lemmas}
\makeatletter
\newlabel{more-algebra-lemma-flat-completion}{{0C4G}{source}{}{sourcepointer-4}{}}
\newlabel{more-algebra-lemma-flat-completion@cref}{{[originallemma][0][]0C4G}{[1][0][]source}{}{}{}}
\makeatother
\title{Formal local structure of a flat family at a split node of the fibre: completed local ring is O\_\{S,s\}\textasciicircum{}wedge[[x, y]]/(xy - h)}
\author{The Stacks Project Authors}
\date{stacks-project:curves.tex:local-snapshot-20261012; GNU FDL 1.2 or later}
\begin{document}
\maketitle
\noindent\textit{Editorial note.} All delivered mathematical text is the author's own, copied from the pinned Stacks Project chapter file curves.tex; no statement or proof step was rewritten, reordered or removed.\par\medskip
\noindent\textit{Editorial note.} The counted proof is the author's complete argument: the localization and completion setup, the split-node presentation of B/mB, the full correction step with its explicit formulas for u', v', h' and delta', the iteration with its Cauchy estimates, and the closing surjectivity and mod-injectivity arguments are all retained verbatim.\par\medskip
\noindent\textit{Editorial note.} The label lemma-formal-local-structure-nodal-family names the counted statement, and the proof environment immediately following it in the source is delivered verbatim as its proof.\par\medskip
\noindent\textit{Editorial note.} The counted statement is carried by the author's own lemma environment in the chapter source.\par\medskip
\noindent\textit{Allowed standard prerequisites (editorial).}\begin{itemize}
\item Split nodes of the central fibre. Same chapter, Lemma lemma-split-node and Definition definition-split-node. Supplies the mod-m normal form: (A/m)[[x, y]] -> B/mB with x, y mapping to u-bar, v-bar is surjective with kernel (xy).
\item Flatness of completed local rings along the family. More on Algebra, Lemma more-algebra-lemma-flat-completion; Algebra, Lemmas algebra-lemma-completion-generalities and algebra-lemma-mod-injective. Gives flatness of A -> B for the completions, surjectivity from the mod-m surjection, and the lifting of mod-m isomorphisms of flat complete algebras.
\end{itemize}\medskip
\ifdefined\hypertarget\hypertarget{sourcepointer-1}{}\fi
\noindent\textit{Source reference.} Original lemma 0315: sources/stacks-chow/tags.txt, source lines 3972--3972. Prerequisite pointer invoked at proof line(s) 5229; it names a standard fact of a sibling chapter that the counted argument uses, and every step of the counted proof built on it is retained verbatim in the delivered body.\par\medskip
\ifdefined\hypertarget\hypertarget{sourcepointer-2}{}\fi
\noindent\textit{Source reference.} Original lemma 00ME: sources/stacks-chow/tags.txt, source lines 820--820. Prerequisite pointer invoked at proof line(s) 5230; it names a standard fact of a sibling chapter that the counted argument uses, and every step of the counted proof built on it is retained verbatim in the delivered body.\par\medskip
\ifdefined\hypertarget\hypertarget{sourcepointer-3}{}\fi
\noindent\textit{Source reference.} Original reference 0CBW: sources/stacks-curves/pinned-original-source.tex, source lines 4585--4585. Prerequisite pointer invoked at proof line(s) 5164; it names a standard fact of the same chapter that the counted argument uses, and every step of the counted proof built on it is retained verbatim in the delivered body.\par\medskip
\ifdefined\hypertarget\hypertarget{sourcepointer-4}{}\fi
\noindent\textit{Source reference.} Original lemma 0C4G: sources/stacks-chow/tags.txt, source lines 15226--15226. Prerequisite pointer invoked at proof line(s) 5160; it names a standard fact of a sibling chapter that the counted argument uses, and every step of the counted proof built on it is retained verbatim in the delivered body.\par\medskip
\begin{lemma}
\label{lemma-formal-local-structure-nodal-family}
Let $f : T \to S$ be a morphism of schemes. Let $t \in T$
with image $s \in S$. Assume
\begin{enumerate}
\item $f$ is flat at $t$,
\item $\mathcal{O}_{S, s}$ is Noetherian,
\item $f$ is locally of finite type,
\item $t$ is a split node of the fibre $T_s$.
\end{enumerate}
Then there exists an $h \in \mathfrak m_s^\wedge$ and an isomorphism
$$
\mathcal{O}_{T, t}^\wedge \cong
\mathcal{O}_{S, s}^\wedge[[x, y]]/(xy - h)
$$
of $\mathcal{O}_{S, s}^\wedge$-algebras.
\end{lemma}

\begin{proof}
We replace $S$ by $\Spec(\mathcal{O}_{S, s})$ and $T$ by the base change
to $\Spec(\mathcal{O}_{S, s})$. Then $T$ is locally Noetherian and hence
$\mathcal{O}_{T, t}$ is Noetherian.
Set $A = \mathcal{O}_{S, s}^\wedge$, $\mathfrak m = \mathfrak m_A$, and
$B = \mathcal{O}_{T, t}^\wedge$. By
More on Algebra, Lemma \ref{more-algebra-lemma-flat-completion}
we see that $A \to B$ is flat. Since
$\mathcal{O}_{T, t}/\mathfrak m_s \mathcal{O}_{T, t} = \mathcal{O}_{T_s, t}$
we see that $B/\mathfrak m B = \mathcal{O}_{T_s, t}^\wedge$.
By assumption (4) and Lemma \ref{lemma-split-node}
we conclude there exist $\overline{u}, \overline{v} \in B/\mathfrak m B$
such that the map
$$
(A/\mathfrak m)[[x, y]] \longrightarrow B/\mathfrak m B,\quad
x \longmapsto \overline{u},
x \longmapsto \overline{v}
$$
is surjective with kernel $(xy)$.

\medskip\noindent
Assume we have $n \geq 1$ and $u, v \in B$ mapping to
$\overline{u}, \overline{v}$ such that
$$
u v = h + \delta
$$
for some $h \in A$ and $\delta \in \mathfrak m^nB$.
We claim that there exist $u', v' \in B$ with
$u - u', v - v' \in \mathfrak m^n B$ such that
$$
u' v' = h' + \delta'
$$
for some $h' \in A$ and $\delta' \in \mathfrak m^{n + 1}B$.
To see this, write $\delta = \sum f_i b_i$ with
$f_i \in \mathfrak m^n$ and $b_i \in B$. Then write
$b_i = a_i + u b_{i, 1} + v b_{i, 2} + \delta_i$ with
$a_i \in A$, $b_{i, 1}, b_{i, 2} \in B$ and $\delta_i \in \mathfrak m B$.
This is possible because the residue field of $B$ agrees with the
residue field of $A$ and the images of $u$ and $v$ in $B/\mathfrak m B$
generate the maximal ideal. Then we set
$$
u' = u - \sum b_{i, 2}f_i,\quad
v' = v - \sum b_{i, 1}f_i
$$
and we obtain
$$
u'v' = h + \delta - \sum (b_{i, 1}u + b_{i, 2}v)f_i + \sum
c_{ij}f_if_j =
h + \sum a_if_i + \sum f_i \delta_i + \sum c_{ij}f_if_j
$$
for some $c_{i, j} \in B$.
Thus we get a formula as above with $h' = h + \sum a_if_i$
and $\delta' = \sum f_i \delta_i + \sum c_{ij}f_if_j$.

\medskip\noindent
Arguing by induction and starting with any lifts $u_1, v_1 \in B$
of $\overline{u}, \overline{v}$ the result of the previous paragraph
shows that we find a sequence of elements
$u_n, v_n \in B$ and $h_n \in A$ such that
$u_n - u_{n + 1} \in \mathfrak m^n B$,
$v_n - v_{n + 1} \in \mathfrak m^n B$,
$h_n - h_{n + 1} \in \mathfrak m^n$,
and such that $u_n v_n - h_n \in \mathfrak m^n B$.
Since $A$ and $B$ are complete we can set
$u_\infty = \lim u_n$, $v_\infty = \lim v_n$,  and
$h_\infty = \lim h_n$, and then we obtain $u_\infty v_\infty = h_\infty$
in $B$. Thus we have an $A$-algebra map
$$
A[[x, y]]/(xy - h_\infty) \longrightarrow B
$$
sending $x$ to $u_\infty$ and $v$ to $v_\infty$.
This is a map of flat $A$-algebras which is an
isomorphism after dividing by $\mathfrak m$.
It is surjective modulo $\mathfrak m$ and hence surjective
by completeness and
Algebra, Lemma \ref{algebra-lemma-completion-generalities}.
Then we can apply Algebra, Lemma \ref{algebra-lemma-mod-injective}
to conclude it is an isomorphism.
\end{proof}

\end{document}
